% GENERATED FILE. Edit canonical lessons, then run npm run generate:books.
% canonical-source-hash: fnv1a64:b2000c642220c86c
% canonical-lessons: ES-C633-rico, ES-C633-miseria, ES-C633-riqueza, ES-C633-lujo, ES-C633-gratuito

\chapter{Ricos y pobres}
\label{ch:es-ricos-y-pobres}
\hlchaptermodality{633}

\begin{chapteropening}
\textbf{By the end of this chapter:} \emph{I can compare rich and poor and say what is free or a luxury.}

Complete the last lesson of chapter 633: I can compare rich and poor and say what is free or a luxury.
\end{chapteropening}

\section[rico]{rico — rich; tasty}

\label{lesson:ES-C633-rico}



\begin{quote}
Before the new one: say the Spanish for to pay in, then the Spanish for the pension.
\end{quote}

\subsection*{You'll want to know: rico}
\textbf{rico} — \textquotedblleft{}rich; tasty\textquotedblright{}.

\begin{quote}
Es un país rico.
\end{quote}

\begin{quote}
\emph{It is a rich country.}
\end{quote}

With \emph{ser} it means rich; with \emph{estar} it means tasty: \emph{la sopa está muy rica}.

\begin{cousinweb}
\emph{rico} was borrowed from Gothic \emph{reiks} 'mighty, powerful', from Germanic \emph{rīkijaz} (EtymDB); being powerful came to mean having money.
\end{cousinweb}

\begin{grammarlens}[title={What you've built}]
One more word for telling, explaining and describing at B1.
\end{grammarlens}

\subsection*{Your turn}
\begin{itemize}
\raggedright
  \item \emph{Say it:} \emph{rico}
  \item \emph{Say it:} \emph{rico}, once more
  \item \emph{Say it:} a sentence of your own with \emph{rico}
  \item \emph{Recall:} say the Spanish for to pay in, then the Spanish for the pension, then say \emph{rico} again
\end{itemize}

\begin{tcolorbox}[breakable,colback=teal!4,colframe=teal!35!black,title={Before you move on}]
What does rico mean? (\textquotedblleft{}Rich; tasty\textquotedblright{}.) Say it once more.
\end{tcolorbox}
\section[la miseria]{la miseria — the poverty, the misery; a pittance}

\label{lesson:ES-C633-miseria}



\begin{quote}
Before the new one: say the Spanish for the pension, then the Spanish for rich.
\end{quote}

\subsection*{You'll want to know: la miseria}
\textbf{la miseria} — \textquotedblleft{}the poverty, the misery; a pittance\textquotedblright{}.

\begin{quote}
Vivían en la miseria.
\end{quote}

\begin{quote}
\emph{They lived in great poverty.}
\end{quote}

\emph{Ganar una miseria} — to earn a pittance.

\begin{cousinweb}
\emph{miseria} is a learned borrowing from Latin \emph{miseria} 'wretchedness' (EtymDB), built on \emph{miser} 'wretched'.
\end{cousinweb}

\begin{grammarlens}[title={What you've built}]
One more word for telling, explaining and describing at B1.
\end{grammarlens}

\subsection*{Your turn}
\begin{itemize}
\raggedright
  \item \emph{Say it:} \emph{la miseria}
  \item \emph{Say it:} \emph{la miseria}, once more
  \item \emph{Say it:} a sentence of your own with \emph{la miseria}
  \item \emph{Recall:} say the Spanish for the pension, then the Spanish for rich, then say \emph{la miseria} again
\end{itemize}

\begin{tcolorbox}[breakable,colback=teal!4,colframe=teal!35!black,title={Before you move on}]
What does la miseria mean? (\textquotedblleft{}The poverty, the misery; a pittance\textquotedblright{}.) Say it once more.
\end{tcolorbox}
\section[la riqueza]{la riqueza — the wealth}

\label{lesson:ES-C633-riqueza}



\begin{quote}
Before the new one: say the Spanish for rich, then the Spanish for the poverty.
\end{quote}

\subsection*{You'll want to know: la riqueza}
\textbf{la riqueza} — \textquotedblleft{}the wealth\textquotedblright{}.

\begin{quote}
La riqueza no lo es todo.
\end{quote}

\begin{quote}
\emph{Wealth is not everything.}
\end{quote}

Built on \emph{rico} with \emph{-eza}.

\begin{cousinweb}
\emph{riqueza} is built on \emph{rico}, which was borrowed from Gothic \emph{reiks} 'mighty' (EtymDB), with the ending \emph{-eza}.
\end{cousinweb}

\begin{grammarlens}[title={What you've built}]
One more word for telling, explaining and describing at B1.
\end{grammarlens}

\subsection*{Your turn}
\begin{itemize}
\raggedright
  \item \emph{Say it:} \emph{la riqueza}
  \item \emph{Say it:} \emph{la riqueza}, once more
  \item \emph{Say it:} a sentence of your own with \emph{la riqueza}
  \item \emph{Recall:} say the Spanish for rich, then the Spanish for the poverty, then say \emph{la riqueza} again
\end{itemize}

\begin{tcolorbox}[breakable,colback=teal!4,colframe=teal!35!black,title={Before you move on}]
What does la riqueza mean? (\textquotedblleft{}The wealth\textquotedblright{}.) Say it once more.
\end{tcolorbox}
\section[el lujo]{el lujo — the luxury}

\label{lesson:ES-C633-lujo}



\begin{quote}
Before the new one: say the Spanish for the poverty, then the Spanish for the wealth.
\end{quote}

\subsection*{You'll want to know: el lujo}
\textbf{el lujo} — \textquotedblleft{}the luxury\textquotedblright{}.

\begin{quote}
Un coche así es un lujo.
\end{quote}

\begin{quote}
\emph{A car like that is a luxury.}
\end{quote}

\emph{De lujo} — luxury, deluxe: \emph{un hotel de lujo}.

\begin{cousinweb}
\emph{lujo} is a learned borrowing from Latin \emph{luxus} 'excess, extravagance' (EtymDB).
\end{cousinweb}

\begin{grammarlens}[title={What you've built}]
One more word for telling, explaining and describing at B1.
\end{grammarlens}

\subsection*{Your turn}
\begin{itemize}
\raggedright
  \item \emph{Say it:} \emph{el lujo}
  \item \emph{Say it:} \emph{el lujo}, once more
  \item \emph{Say it:} a sentence of your own with \emph{el lujo}
  \item \emph{Recall:} say the Spanish for the poverty, then the Spanish for the wealth, then say \emph{el lujo} again
\end{itemize}

\begin{tcolorbox}[breakable,colback=teal!4,colframe=teal!35!black,title={Before you move on}]
What does el lujo mean? (\textquotedblleft{}The luxury\textquotedblright{}.) Say it once more.
\end{tcolorbox}
\section[gratuito]{gratuito — free of charge}

\label{lesson:ES-C633-gratuito}



\begin{quote}
Before the new one: say the Spanish for the wealth, then the Spanish for the luxury.
\end{quote}

\subsection*{You'll want to know: gratuito}
\textbf{gratuito} — \textquotedblleft{}free of charge\textquotedblright{}.

\begin{quote}
La entrada es gratuita.
\end{quote}

\begin{quote}
\emph{Entry is free.}
\end{quote}

More formal than \emph{gratis}.

\begin{cousinweb}
\emph{gratuito} is a learned borrowing from Latin \emph{gratuitus} 'freely given', built on \emph{gratia} 'favour, kindness' (Lewis \& Short).
\end{cousinweb}

\begin{grammarlens}[title={What you've built}]
One more word for telling, explaining and describing at B1.
\end{grammarlens}

\subsection*{Your turn}
\begin{itemize}
\raggedright
  \item \emph{Say it:} \emph{gratuito}
  \item \emph{Say it:} \emph{gratuito}, once more
  \item \emph{Say it:} a sentence of your own with \emph{gratuito}
  \item \emph{Recall:} say the Spanish for the wealth, then the Spanish for the luxury, then say \emph{gratuito} again
\end{itemize}

\begin{tcolorbox}[breakable,colback=teal!4,colframe=teal!35!black,title={Before you move on}]
What does gratuito mean? (\textquotedblleft{}Free of charge\textquotedblright{}.) Say it once more.
\end{tcolorbox}
